% =====================================================================
%  ТАВ. СИСТЕМИЙН ЗАГВАР
% =====================================================================
\chapter{СИСТЕМИЙН ЗАГВАР}

\section{Системийн үйл ажиллагааны урсгал}

Системийн үндсэн урсгал нь NFC уншилтаас эхлээд захиалга хүргэгдэх
хүртэлх процессыг хамардаг. Энэ хүрээнд хэрэглэгч NFC стикерийг
утсаар уншуулснаар ширээ автоматаар танигдаж, цэс нээгдэж, захиалга
өгөх, төлбөр төлөх, захиалгын төлөвийг хянах боломжтой болно.

\begin{figure}[h]
  \centering
  \begin{tikzpicture}[node distance=8mm and 7mm]
    \node[flowbox] (tag) {NFC стикер\\унших};
    \node[flowbox, below=of tag] (menu) {Цахим цэс нээх\\(ширеэ таних)};
    \node[flowbox, below=of menu] (cart) {Захиалга бүрдүүлэх};
    \node[flowbox, below=of cart] (pay) {Төлбөр төлөх\\(QPay/SocialPay/Карт)};
    \node[flowbox, below=of pay] (kitchen) {Салбар хүлээн авах\\бэлтгэх};
    \node[flowbox, below=of kitchen] (deliver) {Ширээнд\\хүргэх};
    \node[st, below=of deliver] (done) {Дуусгав};
    \draw[-{Stealth[length=2.6mm]}] (tag) -- (menu);
    \draw[-{Stealth[length=2.6mm]}] (menu) -- (cart);
    \draw[-{Stealth[length=2.6mm]}] (cart) -- (pay);
    \draw[-{Stealth[length=2.6mm]}] (pay) -- (kitchen);
    \draw[-{Stealth[length=2.6mm]}] (kitchen) -- (deliver);
    \draw[-{Stealth[length=2.6mm]}] (deliver) -- (done);
    \draw[-{Stealth[length=2.6mm]}, bend right=25]
      (deliver.west) to node[left, font=\scriptsize] {Төлбөр төлөгдөөгүй}
      (pay.west);
  \end{tikzpicture}
  \caption{Захиалгын үндсэн урсгалын диаграмм}
  \label{fig:ch5-ursgal}
\end{figure}

\section{Технологийн архитектур}

Систем нь гурван давхаргат архитектуртай \cite{fielding2000, sommerville2016}:

\begin{itemize}
  \item \textbf{Клиент давхарга:} Хэрэглэгчийн PWA (React) болон
        салбарын ажилтны dashboard (React).
  \item \textbf{Аппликешн давхарга:} NestJS сервер талын код,
        REST API, WebSocket бодит цагийн мэдэгдэл.
  \item \textbf{Өгөгдлийн давхарга:} PostgreSQL (үндсэн өгөгдөл),
        Redis (session, queue, cache).
\end{itemize}

\begin{figure}[h]
  \centering
  \begin{tikzpicture}
    \draw[layer] (-0.6,-1.2) rectangle (12.1,1.6);
    \node[layerlbl] at (-0.5,1.5) {Клиент давхарга};
    \node[svc] (pwa) at (2.0,0.2) {Хэрэглэгчийн PWA\\(React)};
    \node[svc] (staff) at (6.0,0.2) {Ажилтны dashboard\\(React)};
    \node[svc] (admin) at (9.8,0.2) {Админ\\хэсэг};

    \draw[layer] (-0.6,-4.6) rectangle (12.1,-2.2);
    \node[layerlbl] at (-0.5,-2.3) {Аппликешн давхарга (NestJS)};
    \node[svc] (api) at (2.0,-3.4) {REST API};
    \node[svc] (ws) at (5.2,-3.4) {WebSocket\\мэдэгдэл};
    \node[svc] (pay) at (8.2,-3.4) {Төлбөрийн модуль};
    \node[ext] (qpay) at (11.0,-3.4) {QPay/SocialPay};

    \draw[layer] (-0.6,-8.0) rectangle (12.1,-5.6);
    \node[layerlbl] at (-0.5,-5.7) {Өгөгдлийн давхарга};
    \node[dbx] (pg) at (2.0,-6.8) {PostgreSQL};
    \node[dbx] (redis) at (5.6,-6.8) {Redis};
    \node[dbx] (s3) at (9.0,-6.8) {Объект хадгаламж\\(зураг файл)};
  \end{tikzpicture}
  \caption{Системийн гурван давхаргат архитектур}
  \label{fig:ch5-arch}
\end{figure}

\section{Өгөгдлийн сангийн загвар}

Системийн үндсэн багцууд: \texttt{User}, \texttt{Venue},
\texttt{Table}, \texttt{MenuCategory}, \texttt{MenuItem}, \texttt{Order},
\texttt{OrderItem}, \texttt{Payment}, \texttt{NfcTag}, \texttt{AuditLog}.

\begin{figure}[h]
  \centering
  \begin{tikzpicture}
    \node[ent] (user) at (0,0) {\texttt{User}\\--- id, name, phone, email};
    \node[ent] (venue) at (5.4,0) {\texttt{Venue}\\--- id, name, address};
    \node[ent] (table) at (10.6,0) {\texttt{Table}\\--- id, venue\_id, no};
    \node[ent] (nfctag) at (10.6,-1.6) {\texttt{NfcTag}\\--- id, table\_id, uid};
    \node[ent] (order) at (5.4,-3.0) {\texttt{Order}\\--- id, table\_id, user\_id, status};
    \node[ent] (orderitem) at (0,-4.4) {\texttt{OrderItem}\\--- id, order\_id, item\_id, qty};
    \node[ent] (menuitem) at (5.4,-5.8) {\texttt{MenuItem}\\--- id, category\_id, name, price};
    \node[ent] (payment) at (10.6,-3.0) {\texttt{Payment}\\--- id, order\_id, method, status};
    \draw[thick] (user) -- (order);
    \draw[thick] (venue) -- (table);
    \draw[thick] (table) -- (nfctag);
    \draw[thick] (table) -- (order);
    \draw[thick] (order) -- (orderitem);
    \draw[thick] (orderitem) -- (menuitem);
    \draw[thick] (order) -- (payment);
  \end{tikzpicture}
  \caption{Өгөгдлийн сангийн объект-харилцааны (ER) диаграмм}
  \label{fig:ch5-er}
\end{figure}

\section{REST API-ийн тодорхойлолт}

\begin{table}[h]
  \centering
  \caption{REST API-ийн үндсэн endpoint-ууд}
  \label{tab:ch5-api}
  \small
  \begin{tabular}{@{}p{2.0cm}p{5.6cm}p{6.4cm}@{}}
    \toprule
    \textbf{HTTP} & \textbf{Endpoint} & \textbf{Тайлбар} \\ \midrule
    GET & \texttt{/api/v1/menu?venue=\{id\}} & Салбарын цэс, ангилал,
          үнийг авах \\
    POST & \texttt{/api/v1/orders} & Шинэ захиалга үүсгэх (ширеэний
           NFC стикерийн uid-ээр) \\
    GET & \texttt{/api/v1/orders/\{id\}} & Захиалгын төлөв, дэлгэрэнгүй
          мэдээлэл \\
    POST & \texttt{/api/v1/payments/intent} & Төлбөрийн санал (intent)
           үүсгэх \\
    POST & \texttt{/api/v1/payments/webhook} & Төлбөрийн үйлчилгээ
           үзүүлэгчийн callback \\
    PATCH & \texttt{/api/v1/staff/orders/\{id\}} & Ажилтан захиалгын
            төлөвийг шилжүүлэх \\
    GET & \texttt{/api/v1/admin/reports} & Орлого, захиалгын тайлан \\
    \bottomrule
  \end{tabular}
\end{table}

\section{Аюулгүй байдлын загвар}

\begin{itemize}
  \item \textbf{Бүх урсгал HTTPS:} TLS 1.2+ протокол, HSTS header
        \cite{rfc9110}.
  \item \textbf{Шошгоны кодчилол:} NFC стикерийн uid-ийг сервер дээр
        баталгаажуулж, шууд бүртгэлээс хамгаалах.
  \item \textbf{Хандалтын эрх:} RBAC (role-based access control) --
        хэрэглэгч, ажилтан, админ гэсэн 3 түвшний эрх
        \cite{owasp}.
  \item \textbf{Төлбөрийн мэдээлэл:} Картын мэдээллийг систем
        хадгалахгүй; төлбөрийн үйлчилгээ үзүүлэгчийн SDK-гаар
        шууд дамжуулна (PCI DSS-ийн шаардлага хангасан \cite{pcidss}).
  \item \textbf{Хувийн мэдээлэл:} «Хувийн нууцын тухай» хуулийн
        дагуу цуглуулсан мэдээллийг зөвхөн үйлчилгээнд шаардлагатай хэмжээнд хадгалж, хугацаа нь болоход устгана
        \cite{privacylaw}.
  \item \textbf{Аудит:} Нэвтрэлт, захиалга, төлбөрийн өөрчлөлтийг
        AuditLog хүснэгтэд хадгалж, гүйлгээний ул мөрийг бүрэн
        мөшгөх \cite{iso27001}.
\end{itemize}

\section{Хөгжүүлэлтийн төлөвлөгөө}

Төслийг 6 сарын хугацаанд 4 үе шаттайгаар хэрэгжүүлнэ
(Хүснэгт~\ref{tab:ch5-tolovlogoo}).

\begin{table}[h]
  \centering
  \caption{Хөгжүүлэлтийн үе шатууд}
  \label{tab:ch5-tolovlogoo}
  \small
  \begin{tabular}{@{}p{1.0cm}p{4.6cm}p{3.6cm}p{4.0cm}@{}}
    \toprule
    \textbf{№} & \textbf{Үе шат} & \textbf{Хугацаа} & \textbf{Гарц (deliverable)} \\
    \midrule
    1 & Шаардлага тодорхойлолт, архитектурын зураглал & 1 -- 2 сар &
    Шаардлагын баримт, диаграммууд \\
    2 & Backend болон өгөгдлийн сангийн хөгжүүлэлт & 2 -- 3 сар &
    REST API, ER бүтэц, төлбөрийн интеграци \\
    3 & PWA болон ажилтны хэсгийн хөгжүүлэлт & 3 -- 5 сар &
    Цэс, захиалга, төлбөрийн интерфэйс \\
    4 & Интеграц, тест, нэвтрүүлэг & 5 -- 6 сар &
    Бүтээгдэхүүний эхний хувилбар (MVP) \\
    \bottomrule
  \end{tabular}
\end{table}

\section{Туршилтын төлөвлөгөө}

\begin{itemize}
  \item \textbf{Нэгж туршилт:} Backend-ийн үндсэн модулиуд (захиалга
        үүсгэх, төлбөрийн intent, төлөвийн шилжилт).
  \item \textbf{Интеграцын туршилт:} NFC → цэс → захиалга → төлбөр →
        ажилтны хүлээн авалт.
  \item \textbf{Хэрэглэгчийн туршилт (UAT):} Анхдагч хэрэглэгч кофе шопт 2 -- 4
        долоо хоногийн турш бодит орчинд турших.
  \item \textbf{Ачааллын туршилт:} Оройн оргил цагт нэг салбарт
        200 -- 500 захиалга/цагийн ачааллыг тэсвэрлэх.
  \item \textbf{Аюулгүй байдлын туршилт:} HTTPS, RBAC, webhook
        баталгаажуулалтын туршилт.
\end{itemize}
